\chapter{Symplectic geometry}

\section{Symplectic manifolds}
\subsection{Lie derivatives}

Some notation that will be used in this subsection:
\begin{center}
    \begin{tabular}{|c|}
        \hline 
        $\Gamma(E)=\Gamma(M,E)$\\
        \hline 
        $\fr{X}(M)=\Gamma(TM)$\\
        \hline 
        $\Omega^k(M)=\Gamma(\bigwedge^kT^*M)$\\
        \hline 
        $\Omega^\bullet(M)=\bigoplus_{k\geq0}\Omega^k(M)$\\
        \hline
    \end{tabular}
\end{center}
Here $E\to M$ is a vecto bundle over $M$.
\begin{dang}
    \begin{definition}
        Let $\pi:E\to M$ be a smooth vector bundle over $M$. Suppose that we have a family of sections 
        \begin{equation}
        \{s_t\}_{t\in I} =  (I\subset\bbR\to \Gamma(M,E):t\mapsto s_t).
        \end{equation}
        We say that $\{s_t\}_{t\in I}$ is smooth if for any $p\in M$, there exist a trivialization 
        $\varphi_U:U\ni p\to U\times\bbR^r$ such that  
        \begin{equation}
            \varphi_U\circ s_t\big|_U=(\mathds{1}_U, a_U(t,\cdot)),\;a_U:I\times U\overset{C^\infty}{\rightarrow}\bbR^r.
        \end{equation}
    \end{definition}
\end{dang}
Of course, this definition does not deppend on the choice of trivialization: If $\varphi_V$ is another trivialization, then 
$a_V(t,\cdot)=g_{VU}a_U(t,\cdot)$, where $\varphi_V\circ\varphi_U^{-1}(p,v)=(p,g_{VU}(p)v)$. Remember that 
$g_{VU}:U\cap V\to GL_r(\bbR)$ is a smooth function. This definition also works for limits. Let $\{s_t\}$ be a 
family of sections of $E$, then we say that $\lim_{t\to t_0}s_t$ exists if $\lim_{t\to t_0}a_U(t,p)$ exists for \textit{all} $p \in U$. 
In this case, we set for $p\in U$
\begin{equation}
    \lim_{t\to t_0}s_t\big|_p=\varphi_U^{-1}(p,\lim_{t\to t_0}a_U(t,p)) \in E.
\end{equation}
Again, this definition does not deppend on the choice of trivialization. In fact 
\begin{equation}
    \begin{split}
        \varphi_V(\lim_{t\to t_0}s_t\big|_p)&=\varphi_V\varphi_U^{-1}(p,\lim_{t\to t_0}a_U(t,p))\\
        &=(p,g_{VU}(p)\cdot\lim_{t\to t_0}a_U(t,p))\\
        &=(p,\lim_{t\to t_0}g_{VU}(p)a_U(t,p))\;(g_{VU}(p)\text{ is linear})\\
        &=(p,\lim_{t\to t_0} a_V(t,p)).
    \end{split}
\end{equation}
The same work for derivatives. For a smooth family $\{s_t\}$ we define 
\begin{equation}
    (\frac{d}{dt}\bigg|_{t=t_0}s_t)_p=\varphi_U^{-1}(p,\frac{\partial}{\partial t}\bigg|_{t=t_0}a_U(t,p))\in E.
\end{equation}
Since linear maps commute with derivatives, by the same calculation made before, we get the nondependence on 
the choice of trivialization.

\begin{dang}
    \begin{remark}
    Since trivialization map are linear maps on the fibers, we have the following relation between limits and derivatives of 
    sections:
    \begin{equation}
        \frac{d}{dt}\bigg|_{t=t_0}s_t=\underset{t\to t_0}{\lim}\frac{s_t-s_{t_0}}{t-t_0}.
    \end{equation}
\end{remark}
\end{dang}

If $\{s_t\}_{t\in[0,1]}$ is a smooth familly of sections of a vector bundle $E$, then we also can define its integral 
$\int_0^1s_tdt \in\Gamma(E)$. Let $(U,\varphi_U)$ be a local trivialization of $E$. If $\varphi_U\circ s_t=(\mathds{1}_U,a_U(t,\cdot))$, then we let 
\begin{equation}
    \int_0^1s_tdt\bigg|_p = \varphi_U^{-1}(p,\int_0^1a_U(t,p)dt),\;\forall p\in U.
\end{equation}
This is well defined because integral of vectors commutes with linear maps, so this definition does not deppend 
on the choice of trivializations. It's clear from this definition that we have 
\begin{equation}
    \int_0^1(\frac{d}{dt}s_t)dt\bigg|_p = s_1(p)-s_0(p),\;\forall p\in M.
\end{equation}
In the case of differential forms, we have good properties. For example, if $\alpha_t,\beta_t\in\Omega^\bullet(M)$, then 
\begin{equation}
    \begin{aligned}
        &\frac{d}{dt}\alpha_t\wedge\beta_t =(\frac{d}{dt}\alpha_t)\wedge\beta_t+\alpha_t\wedge(\frac{d}{dt}\beta_t),\\
        &(d\circ \frac{d}{dt})\alpha_t=(\frac{d}{dt}\circ d)\alpha_t,\\
        &\frac{d}{dt}F^*(\alpha_t)=F^*(\frac{d}{dt}\alpha_t)\quad(\text{comutes with pullback}).
    \end{aligned}
\end{equation}
\begin{remark}
    In the previous equations we mean: 
    \begin{equation}
        \frac{d}{dt}\alpha_t = \frac{d}{ds}\bigg|_{s=t}\alpha_s.
    \end{equation}
\end{remark}


To define Lie derivatives we need to remember the definition of flows of vector fields.
\begin{dang}
    \begin{theorem}
        Let $X\in \fr{X}(M)$ be a smooth vector field on $M$. Then for any $p\in M$, there exist a 
        neighborhood $U\ni p$ and a smooth map called \text{flow} of the vector field $X$, given by 
        \begin{equation}
            \varphi: I\times U\rightarrow M
        \end{equation}
        satisfying
        \begin{equation}
            \frac{\partial}{\partial t}\varphi_t(q)=X_{\varphi_t(q)},\text{ and }\varphi_0=\mathds{1}_U.
        \end{equation}
    \end{theorem}
\end{dang}
So $t\mapsto \varphi_t(q)$ is an integral curve for $X$ with initial point $q$. By uniqueness of integral curves, we can see 
that $\varphi_{t+s}=\varphi_t\circ\varphi_s$ whenever this equation makes sense. So, in particular 
$\varphi_{t^{-1}}=(\varphi_{t})^{-1}$. That is $\varphi_t$ is a diffeomorphism onto its image. With this theorem/definition we 
already can define Lie derivatives.

\begin{dang}
    \begin{definition}
        (Lie derivatives). Let $X,Y\in\fr{X}(M),\Omega\in\Omega^k(M),p\in M$ 
        and $\varphi$ a local flow of $X$ defined in some neighborhood of $p$. We set 
        \begin{equation}
            \begin{aligned}
                &\text{(Vector fields) }\mathcal{L}_X(Y)_p=\frac{d}{dt}\bigg|_{t=0}(\varphi_{-t*}Y)_p\\
                &\text{(Differential forms) }\mathcal{L}_X(\omega)=\frac{d}{dt}\bigg|_{t=0}(\varphi_t^*\omega)_p.
            \end{aligned}
        \end{equation}
    \end{definition}
\end{dang}
A fundamental calculation shows that $\mathcal{L}_X(Y)=[X,Y]$ and if $f\in\Omega^0(M)=C^\infty(M)$, then 
$\mathcal{L}_X(f)=Xf$.\par 

The Lie derivative satisfies the following identities:
\begin{center}
    \begin{tabular}{|c|c|}
        \hline 
        Leibniz rule & $\mathcal{L}_X(\omega\wedge\eta)=\mathcal{L}_X(\omega)\wedge\eta+\omega\wedge\mathcal{L}_X(\eta)$\\
        \hline 
        Commutes with $d$ & $d\circ\mathcal{L}=\mathcal{L}\circ d$\\
        \hline 
        Cartan magic formula & $\mathcal{L}_X=d\circ \iota_X+\iota_X \circ d$\\
        \hline 
        Interior mult. of the commutator & $\iota_{[X,Y]}=\mathcal{L}_X\iota_Y-\iota_Y\mathcal{L}_X$\\
        \hline
    \end{tabular}
\end{center}
So at this point we have the interplay of various endomorphisms of $\Omega^\bullet(M)$:
\begin{center}
    \begin{tabular}{|c|}
        \hline 
       $\mathcal{L}_X:\Omega^\bullet(M)\to\Omega^\bullet(M)$\\
        \hline 
       $d:\Omega^\bullet(M)\to\Omega^{\bullet+1}(M)$\\
        \hline 
       $\iota_X:\Omega^\bullet(M)\to\Omega^{\bullet-1}(M)$\\
       \hline 
    \end{tabular}
\end{center}
Remember the definition of interior multiplication $\iota_X\omega(-)=\omega(X,-)$. It's a nilpotent operator of order 2, 
satisfying the Poincaré rule for wedge multiplication:
\begin{equation}
    \iota_X(\omega\wedge\eta)=\iota_X(\omega)\wedge\eta+(-1)^{k}\omega\wedge\iota_X(\eta),\;\omega\in\Omega^k(M).
\end{equation}

\subsection{Symplectic and Hamiltonian vector fields}
\begin{dang}
    \begin{definition}
        A symplectic manifold is a pair $(M,\omega)$ where $\omega\in\Omega^2(M)$ is a closed and 
        \text{nondenerate} 2-form. By nondegeneracy of $\omega$, we mean that $\omega_p\in\bigwedge^2T_p^*M$ is nondegenerate 
        for all $p\in M$.
    \end{definition}
\end{dang}
Since $\omega_p:T_pM\oplus T_pM\to\bbR$ is nondegenerate, we have an isomorphism of fibers:
\begin{equation}
    T_pM\to T_p^*M: X\mapsto (Y\mapsto \omega_p(X,Y)).
\end{equation}

Let $E,E^\prime$ be two smooth vector bundles over $M$. A vector bundle morphism from $E$ to $E^\prime$ is a smooth map 
$f:E\to E^\prime$, linear on fibers, that makes the following diagram commute:
\begin{center}
    \begin{tikzcd}
        E\arrow[r,"f"]\arrow[d,"\pi"]&E^\prime\arrow[d,"\pi^\prime"]\\
        M\arrow[r,"\mathds{1}_M"]&M
    \end{tikzcd}
\end{center}
The identity morphism from $E$ to $E$ is defined by the map $\mathds{1}_E:E\to E$. Now, we claim that if a morphism of 
vector bundles $f:E\to E^\prime$ is an isomorphism of vector spaces on each fiber, then it's actualy an isomorphism of 
vector bundles. In fact, choose trivializations $(U,\varphi_U)$ and $(U,\varphi_U^\prime)$ of $E$ and $E^\prime$, then 
\begin{equation}
F_U= \varphi_U^\prime\circ f|_U\circ \varphi_U^{-1}:(p,v)\in U\times\bbR^r\mapsto (p,a_U(p)\cdot v)    
\end{equation} 
where $a_U: U\to GL_r(\bbR)$ is a smooth map. Here we are abusing notation and writing $f|_{U}$ in place of $f|_{E|_U}$. 
Therefore, we have a smooth inverse 
\begin{equation}
    F_U^{-1}:U\times\bbR^r\to U\times\bbR^r:(p,v^\prime)\mapsto (p,a_U(p)^{-1}\cdot v^\prime)
\end{equation}
which yields an inverse $g_U=\varphi_U\circ F_U^{-1}\circ (\varphi_U^\prime)^{-1}$ to $f|_U$. In order to define a global 
inverse to $f$, we need to prove that the family $\{g_U\}$ glues to a global map. However, we have on $U\cap V$
\begin{equation}
    \begin{split}
        &g_U\big|_{U\cap V}= g_V\big|_{U\cap V}\iff\\
        &\varphi_U\circ F_U^{-1}\circ (\varphi_U^\prime)^{-1}= \varphi_V\circ F_V^{-1}\circ(\varphi_V^\prime)^{-1}\iff\\
        &F_U = \varphi_U^\prime\circ(\varphi_V^\prime)^{-1}\circ F_V\circ \varphi_V\circ\varphi_U^{-1}\iff \\
        &F_U=\varphi_U^\prime\circ f\big|_{U\cap V}\circ v\varphi_U^{-1}
    \end{split}
\end{equation}
which is true on $U\cap V$. So we gave a global inverse $g:E^\prime\to E$, which is a morphism of vector bundles that 
is the inverse morphism to $f$.\par 

From all this discussion, we see that a nondegenerate 2-form yields a vector bundle isomorphism
\begin{equation}
    \omega^\flat : TM\overset{\cong}{\rightarrow} T^*M.
\end{equation}
Remember the identification
\begin{equation}
    \left\{
        \begin{aligned}
            &\quad\quad\text{Alternating maps}\\
            &\underset{k\text{ times}}{\underbrace{\fr{X}(M)\times\cdots\times \fr{X}(M)}}\to C^\infty(M)
        \end{aligned}
    \right\}\longleftrightarrow
    \Omega^k(M)
\end{equation}
The key point here is the fact that $\omega:M\to \bigwedge^k M$ is smooth if and only if $\omega(X_1,\dots, X_k)$ 
is smooth for any sequence $X_i\in\fr{X}(M),\;i=1,\dots, k$. So, the isomorphism of bundles $\omega^\flat$ yields an isomorphism of sections 
$\fr{X}(M)\to \Omega^1(M):X\mapsto \iota_X(\omega)$. 

\begin{dang}
    \begin{definition}
        (Symplectic and Hamiltonian vector fields).
        Let $(M,\omega)$ be a symplectic manifold.
        A vector field $X\in\fr{X}(M)$ is said to be symplectic if $\mathcal{L}_X(\omega)=0$. Furthermore, $X$ is said 
        to be Hamiltonian if $\iota_X(\omega)=df$ for some $f\in C^\infty(M)$.
    \end{definition}
\end{dang}

By the Cartan magic formula we see that $\iota_X(\omega)=d\iota_X(\omega)$, so $X$ is symplectic if and only if 
$\iota_X\omega$ is closed. Also, by definition we see that $X$ is Hamiltonian if and only if $\iota_X\omega$ is exact.
Then we have a inclusion $\fr{X}_{\text{Ham}}(M)\subset\fr{X}_{\omega}(M)$ and the equality occurs only if 
$H_{\text{dR}}^1(M)=0$.

Remember the followin important facts:
\begin{center}
    \begin{tabular}{|c|c|}
        \hline 
        $i:S\hookrightarrow M$ embbeding, with closed image 
        & $i^*:H_{\text{dR}}^\bullet(M)\to H_{\text{dR}}^\bullet(S)$ surjective\\
        \hline
        $p:M\to N$ surjective submersion
        & $p^*:H_{\text{dR}}^\bullet(N)\to H_{\text{dR}}^\bullet(M)$ injective\\
        \hline
    \end{tabular}
\end{center}


\begin{example}
    (Symplectic vector field on a Torus).
 On the circle $\mathbb{S}^1$, viewed as a regular submanifold of $\bbR^2$, define the closed (but not exact) 1-form 
 \begin{equation}
    \delta\theta = (ydx-xdy)\big|_{\mathbb{S}^1}.
 \end{equation}
 Remember that under the identification $T_p\mathbb{S}^1\cong i_{*,p}T_p\mathbb{S}^1$, we have $T_p\mathbb{S}^1=\Ker f_{*,p}$
 where $f:\bbR^2\to\bbR:(x,y)\mapsto x^2+y^2-1$. So we can see that 
 \begin{equation}
    (y\frac{\partial}{\partial x}-x\frac{\partial}{\partial y})\bigg|_p\in T_p\mathbb{S}^1,\;\forall p\in \mathbb{S}^1.
 \end{equation}
 So there exist $X\in\fr{X}(\mathbb{S}^1)$ such that 
 \begin{equation}
    i_{*,p}X_p=(y\frac{\partial}{\partial x}-x\frac{\partial}{\partial y})\bigg|_p,\;\forall p\in\mathbb{S}^1.
 \end{equation}
 So by definition $\delta\theta(X)=1\in C^\infty(\mathbb{S}^1)$. Now, on the torus $\Sigma = \mathbb{S}^1\times \mathbb{S}^1$, define 
 the form $\omega=\pi_1^*\delta\theta_1\wedge\pi_2^*\delta\theta_2$. Under the identifications
 $T(\mathbb{S}^1\times\mathbb{S}^1)\cong T\mathbb{S}^1\times T\mathbb{S}^1$ and $(X_1,0)\to X_1,\;(0,X_2)\to X_2$ we have 
 the global frame   
 $\{X_1,X_2\}$ of $T(\mathbb{S}^1\times\mathbb{S}^1)$ (in particular we see that $\Sigma$ is paralelizable). 
 In this frame, we compute 
 the matrix of the form $\omega$ at any point. It gives 
 \begin{equation}
   \iota_{X_1}\omega:\left\{\begin{aligned}
        &X_1\mapsto 0\\
        &X_2\mapsto 1
    \end{aligned}\right.
 \end{equation}
 So $\omega$ has matrix representation 
 $\left(\begin{smallmatrix}\;\;0&1\\-1&0\end{smallmatrix}\right)$, which is nondegenerate. 
Then we see that $(\Sigma, \omega)$ is a symplectic manifold. In particular we see that 
\begin{equation}
    \begin{split}
         \iota_{X_1}\omega&=\iota_X(\pi_1^*\delta\theta_1\wedge\pi_2^*\delta\theta_2)\\
         &=\iota_X(\pi_1^*\delta\theta_1)\wedge\pi_2^*\delta\theta_2-\pi_1^*\delta\theta_1\wedge\iota_X(\pi_2^*\delta\theta_2)\\
         &=\iota_X(\pi_1^*\delta\theta_1)\wedge\pi_2^*\delta\theta_2\\
         &=\pi_2^*\delta\theta_2
    \end{split}
\end{equation}
 which is not an exact form. In fact, if $\pi_2^*\delta\theta_2$ is exact, then 
 \begin{equation}
    [\pi_2^*\delta\theta_2]=\pi_2^* [\delta\theta_2]=0\Rightarrow [\delta\theta_2]=0
 \end{equation}
 and $\delta\theta_2$ is exact, which is not the case. Therefore $X_1$ is symplectic, but not 
 Hamiltonian !
\end{example}

\begin{example}
    (Hamiltonian vector fields in $\bbR^{2n}$). Let $\omega=\sum_{i=1}^ndq_i\wedge dp_i$ be the standard 
    symplectic form on $\bbR^{2n}$ where $(q,p)$ is the coordinate chart on $\bbR^{2n}$. Then if $H\in C^\infty(\bbR^{2n})$, 
    we have the corresponding Hamiltonian vector field 
    \begin{equation}
        X_H=\sum_ia_i\frac{\partial}{\partial q_i}+\sum_jb_j\frac{\partial}{\partial p_j}.
    \end{equation}
    Now 
    \begin{equation}
        \begin{split}
             \iota_{X_H}\omega &=\sum_i (\iota_{X_H}dq_i)\wedge dp_i-dq_i\wedge (\iota_{X_H}dp_i)\\
             &=\sum_i (a_idp_i-b_idq_i).
        \end{split}
    \end{equation}
    So $\iota_{X_H}\omega = dH$ if and only if 
    \begin{equation}
        b_{i}=-\frac{\partial H}{\partial q_i}\text{ and }a_i=\frac{\partial H}{\partial p_i},
    \end{equation}
    that is 
    \begin{equation}
        X_H = \sum_i (\frac{\partial H}{\partial p_i}\frac{\partial}{\partial q_i}
        -\frac{\partial H}{\partial q_i}\frac{\partial}{\partial p_j}).
    \end{equation}
\end{example}

\subsection{Examples/counter examples of symplectic manifolds}
\begin{dang}
    \begin{lemma}\label{lemma: Liouville volume}
        (Liouville volume).
        Let $(V, \omega)$ be a vector space with a skew symmetric bilinear form. Then 
\begin{equation}
    (V,\omega)\text{ is symplectic }\iff\;\omega^n\neq 0.
\end{equation}
    \end{lemma}
\end{dang}
Proof: We do the proof in three steps
\begin{itemize}
    \item A necessary condition to $\omega$ be nondegenerate is to $\dim(V)$ be even. If $\{v_1,\dots,v_{2n}\}$ is a 
    basis of $V$ and $\{\alpha_1,\dots,\alpha_{2n}\}$ is the respective dual basis of $V^*$, then we compute the matrix 
    of $\omega^\flat: V\to V^*$ relative to these basis. Suppose that 
    \begin{equation}
        \omega = \sum_{i<j}a_{ij}\alpha_i\wedge\alpha_j
    \end{equation}
    Then we claim that $\omega=(A_{ij})$ where 
    \begin{equation}
       \left\{ \begin{aligned}
            &A_{ij}=a_{ji},\;i>j\\
            &A_{ii}=0,\;\forall i\\
            &A_{ij}=-a_{ij},\;i<j\\
        \end{aligned}\right.
    \end{equation}
    In fact
    \begin{equation}
        \begin{split}
            \omega^\flat(v_k)&=\sum_{i<j}a_{ij}\iota_{v_k}(\alpha_i\wedge\alpha_j)\\
            &=\sum_{i<j}a_{ij}(\iota_{v_k}\alpha_i\wedge\alpha_j-\alpha_i\iota_{v_k}\alpha_j)\\
            &=\sum_{i<j}a_{ij}(\delta_{ik}\alpha_j-\delta_{jk}\alpha_i)\\
            &=\sum_{j=k+1}^{2n}a_{kj}\alpha_j-\sum_{i=1}^{k-1}a_{ik}\alpha_i\\
            &=\sum_{j=k+1}^{2n}A_{jk}\alpha_j+\sum_{i=1}^{k-1}A_{ik}\alpha_i\\
            &=\sum_{i=1}^{2n}A_{ik}\alpha_i.
        \end{split}
    \end{equation}
    \begin{remark}
        We have $\omega(v_k,v_l)=A_{lk}$, so $(A_{ij})$ is the GRAM matrix of $-\omega$.
    \end{remark}
    \item Now, let $\{\beta_1,\dots, \beta_{2n}\}$ be another basis of $V^*$ and $P$ be the change of basis matrix, that is 
    $\alpha_i=\sum_{j}P_{ij}\beta_j$. Then we have 
    \begin{equation}
        \begin{split}
            \omega &= \sum_{i<j}a_{ij}\sum_{k<l}(P_{ik}P_{jl}-P_{il}P_{jk})\beta_k\wedge\beta_l\\
            &=\sum_{k<l}\sum_{i<j}(P_{jl}A_{ji}P_{ik}+P_{il}A_{ij}P_{jk})\beta_k\wedge\beta_l\\
            &=\sum_{k<l}(P^TAP)_{kl}\beta_k\wedge\beta_l.
        \end{split}
    \end{equation}
    So $A$ transforms through $P^TAP$.
    \item Since $A$ is a skew symmetric matrix, we can found a basis $\{\beta_1,\dots,\beta_{2n}\}$ such that 
    \begin{equation}
        P^TAP = \text{Diag}\left(
        \left[\begin{matrix}
            0&-\lambda_1\\\lambda_1&0
        \end{matrix}\right],\dots,
        \left[\begin{matrix}
            0&-\lambda_{n}\\\lambda_n&0
        \end{matrix}\right]
        \right)
    \end{equation}
    with $\lambda_i\in\{0,1\}$. That is
    \begin{equation}
        \omega = \sum_{j=1}^n\lambda_j\beta_{2j-1}\wedge\beta_{2j}.
    \end{equation}
    so we see that $\omega^n=n!(\prod_{j=1}^n\lambda_j)\beta_1\wedge\cdots\wedge\beta_{2n}$. Moreover, if 
    $\{w_1,\dots,w_{2n}\}$ is the relative dual of $\{\beta_i\}$ in $V$, then $\omega^\flat:V\to V^*$ has matrix representation
    $P^TAP$ as before. In particular $\det(\omega)\neq0\iff\det(P^TAP)\neq0\iff \prod_j\lambda_j^2\neq0$, i.e., if 
    $\lambda_j=1$ for all $j$. Therefore  
    \begin{equation}
        \omega^n\neq0\iff \prod_{j=1}^n\lambda_j\neq0\iff \lambda_j=1,\;\forall j\iff \omega\text{ is nondegenerate}.
    \end{equation}
\end{itemize}
\qed

\begin{remark}
    If $\omega\in\bigwedge^2 V$ is a skew symmetric bilinear form, then there exist a basis $v_1,\dots,v_n$ of $V$ such that 
    $\omega$ has matrix representation 
    \begin{equation}
        \text{Diag}\left(
            \left[\begin{matrix}
            0&-1\\1&0
        \end{matrix}\right],\dots,
        \left[\begin{matrix}
            0&-1\\1&0
        \end{matrix}\right],
        [0],\dots,[0]
        \right).
    \end{equation}
    To see this, suppose that $\omega\neq0$, then there exist $u_1,u_2\in V$ such that $\omega(u_1,u_2)\neq0$. We can renormalize
    these vectors in order to get $\omega(u_1,u_2)=1$. Now, let $U=\text{Span}\{u_1,u_2\}$, we claim that 
    \begin{equation}
        V = U\oplus U^\omega,\;U^\omega=\{v\in V:\omega(u,v)=0\;\forall u\in U\}.
    \end{equation}
    First we observe that $U\cap U^\omega=(0)$. In fact, if $u=a_1u_1+a_2u_2\in U^\omega$, then 
    \begin{equation}
        \begin{aligned}
            &B(u,u_1)=a_2u_2 = 0\Rightarrow a_2=0,\\
            &B(u,u_2)=a_1u_1 = 0\Rightarrow a_1=0.
        \end{aligned}
    \end{equation}
    Moreover, if $v\in V$, define $u=\omega(v,u_2)u_1+\omega(v,u_1)u_2\in U$. Then we see that
    \begin{equation}
        \omega(v,u_1)=\omega(u,u_1)\text{ and }\omega(v,u_2)=\omega(u,u_2).
    \end{equation} 
    That is $v-u\in U^\omega$. So we see that $v=u+(v-u)\in U+U^\omega$. An induction step finishes the proof of the claim.
\end{remark}

Now we will see some examples/counter examples of symplectic manifolds. The following lemma is useful.
\begin{dang}
    \begin{lemma}
        Let $(M,\omega)$ be a compact, symplectic manifold. Then $H_{\text{dR}}^2(M)\neq0$.
    \end{lemma}
\end{dang}
Proof: We will show that $[\omega]\in H_{\text{dR}}^2(M)\setminus0$. By definition, we know that 
$\omega_p\in\bigwedge T_p^*M$ is nondegenerate. By lemma~(\ref{lemma: Liouville volume}), we know that 
$\omega_p^n\neq0$, then we see that $\omega^n$ is a volume form and $\int_M\omega^n\neq0$. Now, if 
$[\omega]=0$, then $[\omega]^n=[\omega^n]=0$, i.e., $\omega$ is exact. This contradicts the fact that 
$\int_M\omega^n\neq0$.\qed

By the previous theorem, we see that compact manifolds $M$ such that $H_{\text{dR}}^2(M)=0$ are not symplectic manifolds.
In particular $\mathbb{S}^{2n}$ has no symplectic structure for $n>1$, because it's a compact manifold with 
$H_{\text{dR}}^2(\mathbb{S}^{2n})=0$.
\begin{example}
    Now we list some examples/counter examples of symplectic manifolds.
    \begin{itemize}
        \item $\mathbb{S}^3\times\bbR$ is symplectic. In fact we have a diffeomorphism
        \begin{equation}
           f: \mathbb{S}^3\times\bbR\overset{\cong}{\rightarrow}\bbR^4\setminus0:(v,t)\mapsto e^tv.
        \end{equation}
        However $\bbR^4\setminus0$ is a symplectic manifold with canonical symplectic form 
        $\omega=\sum_{i=1}^2dq_i\wedge dp_i$.
        \item $\mathbb{S}^3\times T^3$ is not symplectic. By applying Kunneth formula we see that 
        \begin{equation}
            H_{\text{dR}}^2(\mathbb{S}^3\times T^3)\cong H_{\text{dR}}^0(\mathbb{S}^3)\otimes H_{\text{dR}}^2(T^3).
        \end{equation}
        Now, suppose that $\omega$ is a symplectic form on $\mathbb{S}^3\times T^3$. Let $a=[\omega]\in H_{\text{dR}}^2(\mathbb{S}^3\times T^3)$. 
        Then, by the previous formula we have 
        \begin{equation}
            a=\sum \alpha_i\otimes \beta_i\Rightarrow a^3=\sum_{i,j,k}\alpha_i\alpha_j\alpha_k\otimes \beta_i\beta_j\beta_k.
        \end{equation} 
        But $\beta_i\beta_j\beta_k\in H_{\text{dR}}^6(T^3)=0$ and $a^3=0$. However, since $\mathbb{S}^3\times T^3$ is compact, 
        the Liouville form $\omega^3$ must be a non exact volume form, which is a contradiction with the equation $a^3=0$.
        \item $\mathbb{S}^3\times \bbR^3$ is symplectic. This follows from a general fact: If $G$ is a Lie group, then 
        we have a vector bundle isomorphism
        \begin{equation}
            T^*G\overset{\cong}{\rightarrow}G\times\fr{g}^*:(g,\beta)\mapsto (g,l_{g,*e}^*\beta).
        \end{equation}
        To see that the previous map is an isomorphism, we observe that it's a vector bundle morphism which is an isomorphism
        of vector spaces on the fibers. To see smoothness, take $g\in G$ and $(U,x_1,\dots, x_n)$ a chart around $g$. Then 
        we have the chain of compositions
        \begin{center}
            \begin{tikzcd}
                U\times\bbR^n\arrow[r]& T^*G\arrow[r]&G\times\fr{g}^*\\
                (g,v)\arrow[r,mapsto] &\sum_iv_idx_i\big|_{g}\arrow[r,mapsto]&(g, \langle v, (l_{g,*e}^*dx_i(g))_{i=1}^n\rangle)
            \end{tikzcd}
        \end{center}
        So we need to check that $g\mapsto l_{g,*e}^*dx_i(g)\in\fr{g}^*$ is smooth. To see this, lets $\gamma: I\to G$ be a curve 
        such that $\gamma(0)=e$, so $\gamma^\prime(0)\in\fr{g}$. Then 
        \begin{equation}
            dx_i\big|_g\circ l_{g,*e}\gamma^\prime(0)=\frac{d}{dt}\bigg|_{t=0}x_i(g\gamma(t))
        \end{equation}
        which is the derivative of the smooth function $(t,g)\mapsto x_i(g\gamma(t))$. To finnish the proof we just 
        observe that $\mathbb{S}^3$ is the Lie group of quaternions.
        \item There is no symplectic, compact submanifold $S\hookrightarrow \bbR^{2n},n\geq1$. Suppose to the countrary: 
        There exist a compact manifold, such that the embbeding $i: S\hookrightarrow \bbR^{2n}$ induces a symplectic structure 
        $i^*(\omega)$ on $S$. Since $i(S)\subset\bbR^{2n}$ is compact, then it's closed. Moreover, since $i$ is an imersion, we 
        see that $i^*:H_{\text{dR}}^2(\bbR^{2n})\to H_{\text{dR}}^2(S)$ is surjective. But 
        \begin{equation}
            H_{\text{dR}}^2(\bbR^{2n})=0\Rightarrow H_{\text{dR}}^2(S)=0.
        \end{equation}
        By the previous lemma, this is a contradiction, because we can not have $H_{\text{dR}}^2=0$ on a compact 
        symplectic manifold.
    \end{itemize}
\end{example}

\section{Symplectomorphisms of the cotangent bundle}
\subsection{Canonical form of the cotangent bundle}
The goal of this section is to show that the cotangent bundle $T^*Q$ of a smooth manifold $Q$ has a 
canonical symplectic structure. In particular, we get a canonical map 
\begin{equation}
    \text{Diff}(Q)\rightarrow \text{Symp}(T^*Q).
\end{equation}

To work with cotangent bundles, it's important to understand in detail how do we define its charts. 
To do this we will set notation. 

\begin{center}
    \begin{tabular}{|c|c|}
        \hline
       $Q$ & a smooth manifold\\
       \hline 
       $M=T^*Q$ & its cotangent bundle\\
       \hline 
       $\pi: M\to Q$ & the canonical projection\\  
        \hline 
        $p$ & a point of $M$\\
        \hline
        $q$ & a point of $Q$ (normaly we will have $q=\pi(q)$)\\
        \hline 
        $\alpha \in \Omega^1(M)$ & the canonical 1-form\\
        \hline
    \end{tabular}
\end{center}

\begin{dang}
    \begin{notation}
        Let $(U,x_1,\dots,x_n)$ be a local chart of $Q$ and let $\bar{x}_i,\xi_i\in C^\infty(TU)\;(i=1,\dots,n)$ be such that 
    \begin{equation}
        \begin{aligned}
            &\bar{x}_i=\pi^*x_i\\
            &\xi_i(p)=p(\frac{\partial}{\partial x_i}\bigg|_q),\;\forall p\in TU.
        \end{aligned}
    \end{equation}
    Then $(TU,\barx_1,\dots,\barx_n,\xi_1,\dots,\xi_n)$ will be a local chart of $M$.
    \end{notation}
\end{dang}

\begin{dang}
    \begin{lemma}
         The cotangent bundle $M$ has a canonical 1-form $\alpha\in\Omega^1(M)$. Let $X\in T_pM$, then we let  
         \begin{equation}
            \alpha_p(X)=  p(\underset{\in T_qQ}{\underbrace{\pi_{*,p}X_p}}),\; p\in T_q^*Q.
         \end{equation}
    \end{lemma}
\end{dang}

We can use the chart defined above to write a local expression to $\alpha$. Let 
\begin{equation}
    \alpha = \sum_ia_id\barx_i+\sum_ib_id\xi_i.
\end{equation}
We have 
\begin{equation}
    \begin{split}
        \alpha(\frac{\partial}{\partial\barx_j})\bigg|_p &=\alpha_p(\frac{\partial}{\partial\barx_j}\bigg|_p)\\
        &=p(\pi_{*,p}(\frac{\partial}{\partial\barx_j}\bigg|_p))\\
        &=p(\frac{\partial}{\partial x_j}\bigg|_q)\\
        &=\xi_i(p)
    \end{split}
\end{equation}
while 
\begin{equation}
    \begin{split}
        \alpha(\xi_i)\bigg|_p &=\alpha_p(\frac{\partial}{\partial\xi_i}\bigg|_p)\\
        &=p(\pi_{*,p}(\frac{\partial}{\partial\xi_i}\bigg|_p))\\
        &=0.
    \end{split}
\end{equation}
So we obtain 
    \begin{equation}
        \alpha\big|_{TU} = \sum_i\xi_id\barx_i.
    \end{equation}
In particular we have  
\begin{equation}
   - d\alpha\big|_{TU} = \sum_i d\barx_i\wedge \xi_i
\end{equation}
which proves that $\omega=-d\alpha$ is nondegenerate 2-form on $M$. It's closed because it's exact, so we have a canonical 
symplectic form defined on $M$. 
\begin{dang}
    \begin{theorem}
        The cotangent bundle $M$ has a symplectic structure given by the form 
        \begin{equation}
            \omega = -d\alpha_{\text{can}}.
        \end{equation}
    \end{theorem}
\end{dang}

The canonical 1-form has an important characterization.
\begin{dang}
    \begin{lemma}
        The canonical 1-form $\alpha\in\Omega^1(M)$ is characterized as being the unique 1-form such that 
        \begin{equation}
            \mu^*\alpha = \mu 
        \end{equation}
        for all $\mu\in \Omega^1(Q)$.
    \end{lemma}
\end{dang}
Proof: First we observe that $\mu: Q\to M$ by definition. So $\mu^*:\Omega^1(M)\to\Omega^1(Q)$ and the equation 
$\mu^*\alpha = \mu$ makes sense. Now, assume that this equation is true and let $p\in M$ with $q=\pi(p)\in Q$. Let 
$(U,x_1,\dots, x_n)$ be a chart centered in $q$ and suppose that $p=\sum a_idx_i(q),a_i\in\bbR$. By defining a new chat 
$y_i=x_i+a_i\;(i=1,\dots,n)$ in $U$ we see that we can assume $x_i(q)=a_i$ without lost of generality. Now 
define a bump function $\rho\in C^\infty(Q)$ with support in $U$ and satisfying 
$\rho|_V\equiv 1$ in some neighborhood $V\ni q,V\subset U$. Also define local sections 
\begin{equation}
    \left\{\begin{aligned}
        &\widetilde{\mu}=\sum x_idx_i\in\Omega^1(U)\Rightarrow \mu=\rho\widetilde{\mu}\in\Omega^1(Q)\\
        &\widetilde{\sigma}=\sum a_idx_i\in\Omega^1(U)\Rightarrow \sigma=\rho\widetilde{\sigma}\in\Omega^1(Q)
    \end{aligned}\right.
\end{equation}
Then we see that $\mu(q)=\sigma(q)=p$. By hypothesis, the following equations are satisfied
\begin{equation}
    \mu^*\alpha=\mu\text{ and }\sigma^*\alpha = \sigma.
\end{equation}
Then we compute
\begin{equation}
    \mu_{*,q}(\frac{\partial}{\partial x_j}\bigg|_q)
    =\frac{\partial}{\partial\barx_j}\bigg|_p+\frac{\partial}{\partial\xi_j}\bigg|_p 
    \text{ and }
    \sigma_{*,q}(\frac{\partial}{\partial x_j}\bigg|_q)=\frac{\partial}{\partial\barx_j}\bigg|_p
\end{equation}
while since $\mu(q)=\sigma(q)=p$, we get 
\begin{equation}
    \mu_q(\frac{\partial}{\partial x_j}\bigg|_q)=\sigma_q(\frac{\partial}{\partial x_j}\bigg|_q)=a_j.
\end{equation}
Therefore, we see that 
\begin{equation}
    \begin{aligned}
        &\alpha_p(\frac{\partial}{\partial\barx_j}\bigg|_p)=a_j=\xi_j(p)\\
        &\alpha_p(\frac{\partial}{\partial\xi_j}\bigg|_p)=0
    \end{aligned}
\end{equation}
so $\alpha = \sum \xi_jd\barx_j$. The converse of the lemma is easy.\qed 

\subsection{Cotangent lifting}
Consider the series of diagrams, where $\phi :Q_1\overset{\cong}{\rightarrow}Q_2$.
\begin{center}
    \begin{tikzcd}
       TQ_1 \arrow[r,"T\phi"]\arrow[d,"\pi"]& TQ_2\arrow[d,"\pi"]\\
       Q_1 \arrow[r,"\phi"]& Q_2
    \end{tikzcd}$\to$
    \begin{tikzcd}
         T^*Q_2 \arrow[r,"T\phi^*"]\arrow[d,"\pi"]& T^*Q_1\arrow[d,"\pi"]\\
       Q_2 \arrow[r,"\phi^{-1}"]& Q_1
    \end{tikzcd}$\to$
    \begin{tikzcd}
        T^*Q_1 \arrow[r,"\hat{\phi}"]\arrow[d,"\pi"]& T^*Q_2\arrow[d,"\pi"]\\
       Q_1 \arrow[r,"\phi"]& Q_2
    \end{tikzcd}
\end{center}
where $\hat{\phi}=(T\phi^*)^{-1}$. We claim that $\hat{\phi}^*\alpha_2=\alpha_1$. This follows from the calculation 
\begin{equation}
    \begin{split}
        (\hat{\phi}^*\alpha_2)_p(X)&=(\alpha_2)_{\hat{\phi}(p)}(\hat{\phi}_{*,p}X)\\
        &=\hat{\phi}(p)\pi_{*,\hat{\phi}(p)}\hat{\phi}_{*,p}X\\
        &=\hat{\phi}(p)(\pi\circ\hat{\phi}_{*,p})X\\
        &=\hat{\phi}(p)(\phi\circ\pi)_{*,p}X\\
        &=\hat{\phi}(p)\phi_{*,\pi(p)}\pi_{*,p}X\\
        &=\hat{\phi}(p)T\phi \pi_{*,p}X\\
        &=((T\phi^*)\hat{\phi}(p))\pi_{*,p}X\\
        &=p(\pi_{*,p}X)\;(\text{because }\hat{\phi}=(T\phi^*)^{-1}).
    \end{split}
\end{equation}

Now, if $Q_1=Q_2=Q$, we have constructed a monomorphism (of groups)  
\begin{equation}
    \text{Diff}(Q)\hookrightarrow \text{Symp}(M):\phi\mapsto \hat{\phi}.
\end{equation}

We will finnish this section with two natural questions: 
\begin{dang}
    \begin{itemize}
    \item[(1)] We know that if $F=\hat{\phi}$ is the cotangent lift of a diffeomorphism $\phi\in\text{Diff}(Q)$, then 
    $F^*\alpha=\alpha$. We will show that the converse is also true, so the equation $F^*\alpha=\alpha$  
    characterizes cotangent lifts of diffeomorphisms.
    \item[(2)] Is the map $ \text{Diff}(Q)\hookrightarrow \text{Symp}(M):\phi\mapsto \hat{\phi}$ surjective? The answer 
    is not. We will show that in general, the following fiber preserving map 
    \begin{equation}
        \phi_A:M\to M: p\mapsto p+A_{\pi(p)}
    \end{equation}
    for $A\in\Omega^1(Q)$ fixed, is not a cotangent lift.
\end{itemize}
\end{dang}
For (1), let $F:M\to M$ be a diffeomorphism such that $F^*\alpha=\alpha$. 
\begin{itemize}
    \item Define the \textit{Euler vector field} $v\in\fr{X}(M)$ as the (unique) solution of the equation
    \begin{equation}
        \iota_v\omega = -\alpha.
    \end{equation}
    In local coordinates we have:
    \begin{equation}
        \iota_v\omega = \sum_i\iota_v(d\barx_i\wedge d\xi_i)=\sum_i d\barx_i(v)d\xi_i-d\xi_i(v)d\barx_i.
    \end{equation}
    So $d\barx_i(v)=0$ and $d\xi_i(v)=\xi_i$, i.e., $v = \sum_i\xi_i\frac{\partial}{\partial\xi_i}$ in local coordinates. 
    Moreover, we claim that $F_*v = v$. To see this we calculate
    \begin{equation}
        \begin{split}
            \iota_{F_*v}\omega(X)\big|_{F(p)}&=\omega_{F(p)}(F_{*,p}v_p,X_{F(p)})\\
            &= (F^*\omega)_p(v_p, (F^{-1})_{*,F(p)}X_{F(p)})\\
            &=\omega_p(v_p, (F^{-1})_{*,F(p)}X_{F(p)})\\
            &=(\iota_v\omega)_p((F^{-1})_{*,F(p)}X_{F(p)})\\
            &=-\alpha_p((F^{-1})_{*,F(p)}X_{F(p)})\\
            &=-((F^{-1})^*\alpha)_{F(p)}X_{F(p)}\\
            &=-\alpha_{F(p)}(X_{F(p)})\\
            &=-\alpha(X)\big|_{F(p)}.
        \end{split}
    \end{equation}
    So, by uniqueness of the solution, we get $F_*v=v$.
    \item Let $\varphi^v$ be the flux of $v$ in some open set $U\subset M$. We claim that $\varphi_t^v\circ F = F\circ\varphi_t^v$. 
    To see this define the curve $\gamma(t)= F^{-1}\circ\varphi_t^v(F(p))$ for some fixed $p\in U$. Then we see that 
    $\gamma(0)= F^{-1}\circ\varphi_0^v(F(p)) = F^{-1}F(p) =p$, and 
    \begin{equation}
       \begin{split}
        \gamma^\prime(t)&= (F^{-1})_{*,\varphi_t^v(F(p))} \frac{d}{dt}\bigg|_t\varphi_t^v(F(p))\\
        &=(F^{-1})_{*,\varphi_t^v(F(p))}v_{\varphi_t^v(F(p))}\\
        &=(F^{-1}_*v)_{F^{-1}\varphi_t^v(F(p))}\\
        &=v_{\gamma(t)}.
       \end{split}
    \end{equation}
    Then we see that $\gamma$ is the integral curve of $v$ starting at $p$. Therefore $\gamma(t)=\varphi_t^v(p)$ by 
    uniqueness of integral curves. Now we compute explictly the flux $\varphi_t^v$ of $v$. We need to solve the EDO 
    \begin{equation}
        \frac{\partial}{\partial t}\varphi_t^v(p) =  v_{\varphi_t^v(p)}.
    \end{equation}
    The left hand side is 
    \begin{equation}
        \sum_i(\barx_i\circ\varphi_t^v(p))^\prime(t)\frac{\partial}{\partial\barx_i}
        +\sum_i(\xi_i\circ\varphi_t^v(p))^\prime(t)\frac{\partial}{\partial\xi_i}.
    \end{equation}
    Using the expression of the Euler vector field obtained before, we get 
    \begin{equation}
        \left\{\begin{aligned}
            &(\barx_i\circ\varphi_t^v(p))^\prime(t)=0\\
            &(\xi_i\circ\varphi_t^v(p))^\prime(t)=\xi_i\circ \varphi_t^v(p)
        \end{aligned}\right.
    \end{equation}
    So $\barx_i\circ\varphi_t^v(p)=\barx_i\circ\varphi_0^v(p)=\barx_i(p)$ and 
    $\xi_i\circ\varphi_t^v(p)=\xi_i(p)e^t$. Therefore $\varphi_t^v(p)=e^tp$.
    \item  We claim that $F(\lambda p)=\lambda F(p)$. 
    By the previous item we have $F(\lambda p)=\lambda F(p)$ for $\lambda>0$. However, for $p\in M$ fixed, the map 
    $f:\lambda\in\bbR\mapsto F(\lambda p)$ is smooth. In fact, by taking limits we see that 
    \begin{equation}
        \lim_{\lambda\to0^+}F(\lambda p)=0_{\pi(F(p))}\text{ and }\lim_{\lambda\to0^-}F(\lambda p)=0_{\pi(F(-p))}
    \end{equation}
    so $\pi(F(p))=\pi(F(-p))$. This shows that we have a smooth map $f:\bbR\to M_{\pi(F(p))}$, 
    of finite dimensional vector spaces. Remember that $M_{\pi(F(p))}\subset M$ is a regular submanifold, so by restricting 
    the codomain we still have a smooth map. Its derivative on 0 yields
    \begin{equation}
        \frac{d}{d\lambda}\bigg|_{\lambda=0}f(\lambda)=\left\{\begin{aligned}
         &\lim_{\lambda\to0^+}\frac{F(\lambda p)}{\lambda}=F(p)\\
         &\lim_{\lambda\to0^-}\frac{F(\lambda p)}{\lambda}=-F(-p)             
        \end{aligned}\right.
    \end{equation}
    So we conclude the claim. Now we build a map $\phi=\pi\circ F\circ s$, where $s\in\Gamma(M)$
    is the zero section, i.e., $s(q)=0_q$. We verify 
    \begin{equation}
        \phi\circ\pi(p)=\pi\circ F(0_{\pi(p)})=\pi(0_{\pi(F(p))})=\pi(F(p))
    \end{equation}
    so $\phi\circ \pi=\pi\circ F$. However, the cotangent lifting also satisfies $\phi\circ\pi=\pi\circ\hat{\phi}$, then we 
    see that $G=\hat{\phi}F^{-1}$ is such that $\pi\circ G=\pi$ while $G^*\alpha=\alpha$. In this condition we claim that 
    $G=\mathds{1}_M$. In fact, we have 
    \begin{equation}
        \begin{split}
            (G^*\alpha)_p(X)&=\alpha_{G(p)}(G_{*,p}X)\\
            &=G(p)\pi_{*,G(p)}G_{*,p}X\\
            &=G(p)(\pi\circ G)_{*,p}X\\
            &=G(p)\pi_{*,p}X
        \end{split}
    \end{equation}
    while $\alpha_p(X)=p(\pi_{*,p}X)$. Since $\pi_{*,p}$ is surjective, we get $G(p)=p$. This finnish the proof.\qed
\end{itemize}

For (2), we calculate 
\begin{equation}
    \begin{split}
        (\phi_A^*\alpha)(X)\bigg|_p&=\alpha_{\phi_A(p)}(\phi_{A*}X_p)\\
        &=\phi_A(p)\pi_{*,\phi_A(p)}\phi_{A*}X_p\\
        &=\phi_A(p)(\pi\circ\phi_A)_{*,p}X_p\\
        &=\phi_A(p)\pi_{*,p}X_p \\
        &=(p+A_{\pi(p)})\pi_{*,p}X_p\\
        &=\alpha_p(X_p)+(\pi^*A)_p(X_p)
    \end{split}
\end{equation}
and we conclude that $\phi_A^*\alpha=\alpha+\pi^*A$. So $\phi_A^*$ is a symplectomorphism of $M$ if and only 
if $A\in\Omega^1(Q)$ is closed. However, by the previous item, we see that if $A\neq0$, then $\phi_A$ is not 
a cotangent lift. 

\section{Poisson Bracket and Darboux coordinates}
\subsection{Poisson Bracket}
Let $\omega\in\Omega^1(M)$ be a nondegenerate form (which is not necesseraly closed). Define the 
bilinear map $C^\infty(M)\times C^\infty(M)\to C^\infty(M)$ by 
\begin{equation}
    \{f,g\}=\omega(X_g,X_f),\; f,g\in C^\infty(M).
\end{equation}
We claim that $\{\;,\;\}$ satisfies the Jacobi identity if and only if $d\omega=0$. In fact, if $f,g,h\in C^\infty(M)$, then 
we have 
\begin{equation}
 -\iota_{[X_f,X_g]}\iota_{X_h}\omega = \iota_{X_h}\iota_{[X_f,X_g]}\omega.
\end{equation}
The LHS is equal to 
\begin{equation}
    \begin{split}
       -\iota_{X_h}\omega([X_f,X_g])&=-dh([X_f,X_g])\\
        &=-(X_fX_g)h+(X_gX_f)h\\
        &=-X_f(X_gh)+X_g(X_fh)\\
        &=-X_f\{g,h\}+X_g\{f,h\}\\
        &=-\{f,\{g,h\}\}+\{g,\{f,h\}\}\\
        &=-\{f,\{g,h\}\}-\{g,\{h,f\}\}.
    \end{split}
\end{equation}
For the RHS we develop first 
\begin{equation}
    \begin{split}
        \iota_{[X_f,X_g]}\omega&=(\mathcal{L}_{X_f}\iota_{X_g}-\iota_{X_g}\mathcal{L}_{X_f})\omega\\
        &=\mathcal{L}_{X_f}dg -\iota_{X_g}\iota_{X_f}d\omega \\
        &=d(\iota_{X_f}dg)-\iota_{X_g}\iota_{X_f}d\omega\\
        &=d(X_fg)-\iota_{X_g}\iota_{X_f}d\omega\\
        &=d\{f,g\}-\iota_{X_g}\iota_{X_f}d\omega.
    \end{split}
\end{equation}
Then we get 
\begin{equation}
    \begin{split}
       \iota_{X_h}d\{f,g\}-\iota_{X_h}\iota_{X_g}\iota_{X_f}d\omega&=\mathcal{L}_{X_h}\{f,g\}-d\omega(X_f,X_g,X_h)\\
       &=\{h,\{f,g\}\}-d\omega(X_f,X_g,X_h)
    \end{split}
\end{equation}
Because $\iota_{X_h}\{f,g\}=0$. Then we get 
\begin{equation}
    J(f,g,h)=d\omega(X_f,X_g,X_h),\;\forall g,f,h\in C^\infty(M)
\end{equation}
where $J(f,g,h)=\{f,\{g,h\}\}+\{g,\{h,f\}\}+\{h,\{f,g\}\}$ is the Jacobiator of $f,g$ and $h$. So the Jacobi identity
is equivalent to $J(f,g,h)=0,\;\forall f,g,h\in C^\infty(M)$. For any point $p$, we have an isomorphism 
$T_pM\overset{\cong}{\to}T_p^*M:X\mapsto\iota_X\omega_p$. But if $(U,x_1,\dots,x_n)$ is a chart, then 
$dx_i(p),i=1,\dots,n$ is a basis of $T_p^*M$ for any $p\in U$. Fix a point $p\in U$ and let $\rho\in C^\infty(M)$ 
be a bump function with support in $U$ satisfying $\rho|_V\equiv 1$ for some open set $V\subset U$ containing $p$. Therefore 
$\widetilde{x}_i=\rho x_i\in C^\infty(M)$ and we have Hamiltonian vector fields $X_i=X_{\widetilde{x}_i},i=1,\dots,n$ 
such that $\text{Span}(X_i(q))_{i=1}^n=T_qM,\;\forall q\in V$. This proves that if $d\omega(X_f,X_g,X_h)=0$ for 
all $f,g,h\in C^\infty(M)$, then $d\omega=0$.\par 

In conclusion, for a symplectic manifold $(M,\omega)$, we have a Lie algebra $(C^\infty(M),\{\cdot,\cdot\})$. Moreover, 
we have the identity $\{f,gh\}=\{f,g\}h+g\{f,h\}$. In fact 
\begin{equation}
    \{f,gh\}=X_f(gh)=(X_fg)h+g(X_fh)=\{f,g\}h+g\{f,h\}.
\end{equation}
Therefore we have an induced Poisson manifold structure on $M$.

\section{Distributions and symplectic reduction}

A distribution on a smooth manifold $M$ is a map $M\ni p\mapsto D_p\subset TpM$. The rank of a distribution 
at a point $p\in M$ is $\dim(D_p)$. Also, we say that $D$ is a smooth distribution if for each point $p\in M$ and 
$v\in D_p$, there exist $X\in\fr{X}(M)$ subordinated to $D$, i.e., $X_q\in D_q$, satisfying $X_p=v$. We have 
the following observation:
\begin{center}
    $D$ is smooth of constant rank $\iff$ $D$ is a subbundle of $TM$.
\end{center}
In fact, suppose that $D$ is a smooth distribution of rank $k$ at $p\in M$. Let $\{v_1,\dots, v_k\}$ a basis 
of $D_p$ and $X_i\in\fr{X}(M)$ the associated vector fields subordinated to $D$ and satisfying $X_i(p)=v_i,\;i=1,\dots,n$. 
Choose a chart $(U,x_1,\dots, x_n)$ around $p$ and let 
\begin{equation}
    X_i=\sum_j a_{ij}\frac{\partial}{\partial x_j},\; a_{ij}\in C^\infty(M).
\end{equation} 
If $U\to  M_{k\times n}(\bbR)$, then $A$ has maximal rank at $p$ and we can permute the $x_j$'s in such a way 
that $A=(a_{ij})_{i,j=1}^k$ has full rank. So $\det(A)\neq0$ in a neighborhood $V\ni p$. So, we see that 
$TV=\text{Span}(X_1,\dots,X_k)\oplus \text{Span}(\tfrac{\partial}{\partial x_{k+1}},\dots,\tfrac{\partial}{\partial x_n})$ 
componentwise. This frame defines $D|_V$ as a regular submanifold of $TV$ of dimension $k$. It's local trivialization is 
defined by the frame $\{X_1,\dots, X_k\}$ at $V$.\par 

We say that a smooth distribution is integrable if for each $p\in M$, there exist a regular submanifold $S\ni p$ such that 
$D|_S=TS$. For each $p$, we order by inclusion these submanifolds and call the maximal ones by \textit{leafs}.
Now we write $\Gamma(D)$ for the set of vector fields of $M$ subordinated to $D$. By the Frobenius theorem, we have
\begin{center}
    $D$ is integrable $\iff$ $[\Gamma(D),\Gamma(D)]\subset \Gamma(D)$.
\end{center} 

Two facts:
\begin{itemize}
    \item If $\omega\in\Omega^2(M)$, then 
    \begin{center}
        $D=\Ker(\omega^\flat)$ is smooth $\iff$ if it's of costant rank.\;\;(why??)
    \end{center}
    \item If $d\omega=0$, then $D$ is involutive. In fact, let $X,Y\in\Gamma(D)$, then 
    \begin{equation}
        \begin{split}
            \iota_{[X,Y]}\omega&=\mathcal{L}_X\iota_Y\omega-\iota_Y\mathcal{L}_X\omega\\
            &=-\iota_Y\iota_Xd\omega\quad(\text{because }\iota_X\omega=\iota_Y\omega=0)\\
            &=0\quad(\text{because }d\omega=0),
        \end{split}
    \end{equation}
    so $[X,Y]\in D$. Then we get the important observation:
    \begin{dang}
        \begin{remark}
        $\text{Ker}(\omega)$ has constant rank if and only if it's an integrable distribution, so its define leafs 
        on $M$.
    \end{remark}
    \end{dang}
\end{itemize}

Before, the next example recall that if $f: M\to N$ is smooth with regular value $c$. Then  
$S=f^{-1}(c)$ is a regular submanifold. Furthermore we have 
\begin{center}
    $j_{*,p}T_pS=\Ker(f_{*,p}),\;\forall p\in S$
\end{center}
where $j:S\to M$ is the inclusion (which is an embedding). So a vector field $X\in\fr{X}(M)$ is tangent to $S$ if and only if 
$f_{*,p}X_p=0$ for all $p\in M$. In particular if $f=(f_1,\dots,f_n):M\to\bbR^n$, then 
$df(X)|_p=(df_i(X)|_p)_{i=1}^n$. That is, $X$ is tangent to $S$ if and only if $df_i(X)=0,i=1,\dots,n$. 

\begin{example}
    Let $M$ be a manifold and $\omega \in \Omega^k(M)$. 
    Suppose that $\pi : M \to B$ is a surjective submersion with connected fibers. 
    We say that $\omega$ is \textit{basic} (with respect to $\pi$) if there exists a form $\overline{\omega} \in \Omega^k(B)$ 
    such that $\pi^*\overline{\omega} = \omega$.

\begin{enumerate}
    \item[(a)] Show that $\omega$ is basic iff $i_X\omega = 0$ and $\mathcal{L}_X\omega = 0$ for all vector fields $X$ tangent 
    to the fibers of $\pi$. In particular, if $\omega$ is closed, show that it is basic if 
    $\ker(T\pi) \subseteq \ker(\omega)$ (pointwise in $M$).\par 
    Solution: $(\Rightarrow)$ Let $\omega=\pi^*\overline{\omega}$ be basic and let $X$ be a vector field tangent to $S=\pi^{-1}(q)$. 
    We have 
    \begin{equation}
        \iota_X\omega(Y)\big|_p=\omega_p(X_p,Y_p)=\overline{\omega}_q(\pi_{*,p}X_p,\pi_{*,p}Y_p)=0
    \end{equation}
    because $\pi_{*,p}X_p=0$. Furthermore, $\omega$ basic implies that $d\omega$ is basic. So $\iota_Xd\omega=0$ and 
    $\mathcal{L}_X\omega=0$.\par 
    $(\Leftarrow)$ Fix $q\in B$ and let $p\in \pi^{-1}(q)$. Since $\pi$ is a submersion, there exist charts 
    $(U,\varphi)$ and $(V,\psi)$ centered around $p$ and $q$ respectively, satisfying $\pi(U)\subset V$ and 
    \begin{equation}
        \psi\circ\pi\circ\varphi^{-1}(r_1,\dots,r_m,r_{m+1},\dots,r_n)=(r_1,\dots,r_m).
    \end{equation}
    Therefore $\pi^{-1}(q)\cap U=\{x_1=0,\dots,x_m=0\}$ where we write $x_i=r_i\circ\varphi$. In particular, we see that 
    $X$ is tangent to $\pi^{-1}(q)$ in $U$ if and only if $dx_i(X)=0,i=1,\dots,m$, i.e., 
    $X\in\text{Span}(\tfrac{\partial}{\partial x_{m+1}},\dots,\tfrac{\partial}{\partial x_n})$ in $U$. Now, in $U$ we have 
    $\omega=\sum_I\omega_Idx_I$ where $I=(i_1,\dots, i_k)$ are increasing sequences of integers between 1 and $n$. The condition 
    $\iota_X\omega=0$ for all vector fields tangent to $\pi^{-1}(q)$, implies that $\omega_I=0$ whenever 
    $I=(i_1,\dots, i_k)$ and $i_\ell\geq m+1$ for some $1\leq\ell\leq k$. We also see that  
    \begin{equation}
        \mathcal{L}_X\omega = \iota_Xd\omega 
        = \iota_X(\sum_{I,j}\frac{\partial\omega_I}{\partial x_j}dx_j\wedge dx_I)=0.
    \end{equation}
    So $\tfrac{\partial \omega_I}{\partial x_j}=0,\forall j\geq m+1$. Assuming $U$ connected, this means that 
    \begin{equation}
        \omega_I\circ\varphi^{-1}(r_1,\dots,r_m,r_{m+1},\dots, r_n)=\omega_I\circ\varphi^{-1}(r_1,\dots,r_m,0,\dots,0)
    \end{equation}
    which means that $\omega_I$ is constant on $\pi^{-1}(q)$. This proves that for all $I=(i_1,\dots, i_k)$, 
    the the map $p\mapsto \omega_p(\partial_{i_1},\dots,\partial_{i_k})$ is constant 
    on $\pi^{-1}(q)$. Since the fibers are connected, 
    this shows that the map $p\mapsto \omega_p(X_1(p),\dots, X_k(p))$ is constant on the fibers 
    for any sequence of vector fields $X_1,\dots, X_k\in\fr{X}(M)$. Now, let $Y_i\in\fr{X}(B),i=1,\dots,k$ be fixed. We want 
    to define $\overline{\omega}(Y_1,\dots,Y_k)$. To do this we observe that since $\pi$ is a surjective submerion, there exist 
    $X_i\in\fr{X}(M),i=1,\dots,k$ such that $\pi_{*,p}X_p=Y_{\pi(p)}$. Now, since the 
    map $\omega(X_1,\dots, X_k)\in C^\infty(M)$ is costant on the fibers of $\pi$, we have an unique smooth map 
    $F\in C^\infty(B)$ such that $F\circ\pi = \omega(X_1,\dots, X_k)$. We then define $\overline{\omega}(Y_1,\dots,Y_k)=F$. Since 
    $\iota_X\omega=0$ whenever $X\in\Ker(\pi_{*})$, we see that this defintion is independent on the choice of vector fields 
    $X_i\in\fr{X}(M)$ satisfying $\pi_{*,p}X_p=Y_{\pi(p)}$.


    \item[(b)] Suppose that $\omega$ is a closed $2$-form on $M$ and $\ker(T\pi) = \ker(\omega)$. 
    Show that $\omega = \pi^*\overline{\omega}$ and $\overline{\omega} \in \Omega^2(B)$ is symplectic.\par 
    Solution: First we need to show that $\omega$ is basic. The vector fields $X\in\fr{X}(M)$ 
    tangent to the fibers of $\pi$ are such that $X_p\in\Ker(\pi_{*,p})$, i.e., $X\in\Gamma(\Ker(\pi_*))$. So the 
    equality $\Gamma(\Ker(\pi_*))=\Gamma(\ker(\omega))$ means that $\iota_X\omega=0$ for all vector fields $X$ which are 
    tangent to the fibers of $\pi$. Moreover, since $\omega$ is closed we see that 
    \begin{equation}
        \mathcal{L}_X\omega = (d\iota_X+\iota_Xd)\omega=0
    \end{equation} 
    as well. By the previous item $\omega$ is basic, i.e., there exist $\overline{\omega}\in\Omega^2(B)$ such that 
    $\pi^*(\overline{\omega})=\omega$. This is a closed form by the injectivity of $\pi^*$ and it's nondegenerate because 
    $\iota_Y\overline{\omega}=0$ if and only if $\iota_X\omega=0$ where $X$ and $Y$ are $\pi$ related. However, the last 
    equality means that $X\in\Gamma(\Ker(\omega))=\Gamma(\Ker(\pi_*))$, so $\pi_*X=0\Rightarrow Y=0$.

    \item[(c)] (Application to reduction) Let $(M,\omega)$ be a symplectic manifold and $i : N \hookrightarrow M$ a 
    submanifold such that $D = TN \cap TN^\omega \subset TN$ has constant rank (e.g., $N$ could be coisotropic). 
    We saw in class that $D$ is an integrable distribution (by Frobenius); suppose that the leafspace $B := N/\sim$ 
    is smooth so that the natural projection $\pi : N \to B$ is a submersion. Show that $B$ inherits a unique 
    symplectic form $\omega_{\text{red}}$ with the property that $\pi^*\omega_{\text{red}} = \iota^*\omega$.\par
    Solution: We know that $D=TN\cap TN^\omega=\Ker(\omega_N)$ where $\omega_N=i^*\omega$. We know that $D$ is of constant rank 
    if and only if it's integrable, so it induces leafs on $N$. If the leaf space $B=N/\sim$ is a smooth manifold, then we have 
    an induced surjective submersion $\pi:N\to B$ such that its fibers are the leafs induced by $D$ on $N$. Now we claim that 
    $\Gamma(D)=\Gamma(\Ker(\pi_*))$. In fact, if $X\in\Gamma(\Ker(\pi_*))$, then $X$ is a vector field tangent to the fibers of $\pi$, 
    which are leafs of $N$ induced by $D$. If $p\in M$, then there exist a leaf $S\ni p$ such that $TS=D|_S$, so 
    $X_p\in T_pS=D_p$, i.e. $X\in\Gamma(D)$. This proves that $D=\Ker(\pi_*)$ and by the previous item we see that 
    there exist $\omega_{\text{red}}\in\Omega^2(B)$ such that $\pi^*\omega_{\text{red}}=\omega_N$.
\end{enumerate}
\end{example}

Now we will apply this machinery to a concrete example.
\begin{example}
    Let $M$ be a symplectic manifold and $H\in C^\infty(M)$ a Hamiltonian. Suppose that $c\in\bbR$ is a regular 
    value of $H$, so $N=H^{-1}(c)\subset M$ is a regular submanifold. Since $\text{codim}(N)=1$, it's a coisotropic 
    submanifold, i.e., $TN^\omega\subset TN$. This follows from the following. If $(V,\Omega)$ is a symplectic vector space 
    and $W\subset V$ is a codimension 1 vector space, then $W^\Omega\subset W$. In fact, since $L=W^\Omega$ is a line, so 
    $L\subset W\iff \exists \ell\in L\cap W$. If this is not the case, then $V=W\oplus L$ and any $v\in V$ is of the form 
    $v=w+\lambda\ell$. So $\Omega(v,\ell)=\Omega(w,\ell)=0$, because $\ell\in L=W^\Omega$ and $\Omega$ is not nondegenerate.\par 

    Now we know that $D=TN\cap TN^\omega=TN^\omega$, which is of constant rank 1. Since $D=\Ker(\omega_N)$, it's an 
    integrable distribution which induces a set of foliations on $N$ called the ``null foliation''. In this case, we 
    have an explictly description of the vector fields subordinated to $D$: $\Gamma(D)=\langle X_H \rangle$. 
    In fact, if $Y\in \fr{X}(N)=\Gamma(\Ker(dH))$, then 
    \begin{equation}
        \omega(X_H,Y)=dH(Y)=0\Rightarrow X_H\in\Gamma(TN^\omega).
    \end{equation}

    Now, for a concrete example, let $H:\bbR^{2n}\to\bbR$ be defined by $H(q,p)=\frac12\sum_i(q_i^2+p_i^2)$ and $c=\frac12$. 
    So we get $\mathbb{S}^{2n-1}=H^{-1}(c)$, which is a coisotropic submanifold which has an induced set of foliations. 
    Locally, these foliations are integral curves of 
    \begin{equation}
        X_H=\sum_i(\frac{\partial H}{\partial p_i}\frac{\partial}{\partial q_i}
        -\frac{\partial H}{\partial q_i}\frac{\partial}{\partial p_i})
        =\sum_i(p_i \frac{\partial}{\partial q_i}  - q_i\frac{\partial}{\partial p_i}).
    \end{equation}
    Therefore, to find the leafs, we solve the ODE
    \begin{equation}
        X_H\circ\gamma(t)=\gamma^\prime(t)=\sum_i(q_i\circ\gamma)^\prime(t)\frac{\partial}{\partial q_i}+
        \sum_i(p_i\circ\gamma)^\prime(t)\frac{\partial}{\partial p_i} \quad \bigg|_{\gamma(t)}
    \end{equation}
    which give us the set of equations
    \begin{equation}
        \left\{\begin{aligned}
            &(q_i\circ\gamma)^\prime(t)=p_i\circ\gamma(t)\\
            &(p_i\circ \gamma)^\prime(t)=-q_i\circ\gamma(t)
        \end{aligned}\right.
    \end{equation}
    Define $z_j(t)=q_j(t)+ip_j(t)$. Then the above set is equivalent to $z_j^\prime(t)=-iz_j(t)$ whose solution is 
    $z_j(t)=e^{-it}z_j(0)$. So if an integral curve of $X_H$ start at $p\in\mathbb{S}^{2n-1}$, then all elements 
    $\lambda\cdot p,\;\lambda\in\mathbb{S}^1$ are on the same leaf. In particular, we see that the induced leafspace is 
    \begin{equation}
        B = \frac{\mathbb{S}^{2n-1}}{\mathbb{S}^1}\cong\bbC P^{n-1}
    \end{equation}
    which is a smooth manifold. By the previous discussion on symplectic reduction, we 
    conclude that $\bbC P^{n-1}$ is a symplectic manifold.
\end{example}

\section{Smooth actions of Lie groups on symplectic manifolds}
\subsection{Review on Lie groups and infinitesimal actions}
\begin{itemize}
    \item Let $F:M\rightarrow N$ be a diffeomorphism and $X\in\fr{X}(M)$. We define the pushforward 
    $F_*X\in\fr{X}(N)$ to be vector field defined by the equation $F_*X |_{F(p)}=F_*X_p$. It's a general fact that 
    if $X,Y\in\fr{X}(M)$, then $F_*[X,Y]=[F_*X,F_*Y]$.\par 
    Moreover, if $X\in\fr{X}(M)$ and $Y\in\fr{X}(N)$, we say that $X$ and $Y$ are $F-$related (and we will write 
    $X\sim_FY$) if $F_*X=Y$.
    \item Let $G$ be a Lie group. Define $\fr{X}_{\text{inv}}(G)=\{X\in\fr{X}(G):\ell_{g*}X=X\}$ to be the set of 
    left invariant vector fields of $G$. Then, it's clear that 
    $[\fr{X}_{\text{inv}}(G),\fr{X}_{\text{inv}}(G)]\subset \fr{X}_{\text{inv}}(G)$, so $\fr{X}_{\text{inv}}(G)$ is a 
    Lie algebra. We have an isomorphism 
    \begin{equation}
        \fr{g}\overset{\cong}{\to}\fr{X}_{\text{inv}}(G)=T_eG: X\mapsto [g\mapsto \ell_{g,*e}X]
    \end{equation}
    As a matter of notation, we will write $u^L$ for the left invariant vector field generated by $u\in\fr{g}$, i.e.,
    $u^L(g)=\ell_{g,*}u$.
    \item A $G$- manifold $M$ is a triple $(M,G,\psi)$ where $G$ is a Lie group and 
    $\psi:G\times M\to M$ is a smooth map such that $\psi_g=\psi(g,\cdot)$ defines a homomorphism 
    $G\to\text{Diff}(M):g\mapsto \psi_g$. An infinitesimal generator of $\psi$ is a vector field defined by 
    \begin{equation}
        u_M(x)=\frac{d}{dt}\bigg|_{t=0}\psi_{\exp(tu)}(X).
    \end{equation}
    This defines a map $\fr{g}\to\fr{X}(M):u\mapsto u_M$ which we will show to be a Lie algebra anti-homomorphism.
    \begin{example}
        (infinitesimal generator of $\Ad$). Consider the adjoint action $G\times\fr{g}\to\fr{g}$. We have 
        \begin{equation}
            \begin{split}
                u_{\fr{g}}(v)&=\frac{d}{dt}\bigg|_{t=0}\Ad_{\exp(tu)}(v)\\
                &=\frac{d}{dt}\bigg|_{t=0}(r_{\exp(-tu)}\circ \ell_{\exp(tu)})_{*,e}(v)\\
                &=\frac{d}{dt}\bigg|_{t=0}d_{\exp(tu)}r_{\exp(-tu)}(v^L(\exp(tu)))\\
                &=\frac{d}{dt}\bigg|_{t=0} r_{\exp(tu)*}(v^L)\big|_e\\
                &=\mathcal{L}_{u^L}(v^L)\big|_e\\
                &=[u^L,v^L]_e\\
                &=[u,v].
            \end{split}
        \end{equation}
        Therefore $\Ad_{*,e}=\ad$.
    \end{example}
    \item We have three important results that we will need: 
        \begin{itemize}
            \item[1.] Let $\varphi:M_1\to M_2$ be a $G$ equivariant map, i.e., $\varphi(gx)=g\varphi(x),\;\forall x\in M_1$, then 
            for any $u\in\fr{g}$, we have $u_{M_1}\sim_\varphi u_{M_2}$.
            \item[2.] Let $M$ be a $G$-manifold, then $\fr{g}\to\fr{X}(M):u\mapsto u_M$ is a Lie algebra anti-homomorphism. 
            In fact, the result is true for $(G,\psi(g,\cdot)=r_{g^{-1}})$. This follows frmo the equation 
            \begin{equation}
                u_G(g)=\frac{d}{dt}\bigg|_{t=0}r_{\exp(-tu)}(g)=-u^L(g).
            \end{equation}
            Remember that $r_{\exp(tu)}$ defines the flux of the left invariant vector field $u^L$. Now, we observe that 
            the result is also true for $\overline{M}=G\times M$ with action $g\cdot(a,x)=(r_{g^{-1}}(a),x)$. In fact 
            \begin{equation}
                \begin{split}
                    \pi_{G*(a,x)}u_{\overline{M}}(a,x)&=\frac{d}{dt}\bigg|_{t=0}\pi_G\circ(r_{\exp(-tu)}(a),x)\\
                    &=\frac{d}{dt}\bigg|_{t=0}r_{\exp(-tu)}(a)\\
                    &=-u^L(a).
                \end{split}
            \end{equation}
            Therefore
            \begin{equation}
                [u_{\overline{M}},v_{\overline{M}}]=
                -([u^L,v^L],0)=
                -([u^L,v^L],0)=
                -[u,v]_{\overline{M}}.
            \end{equation}
            Now define the $G$-equivariant map $F:\overline{M}\to M:(a,x)\mapsto a^{-1}x$. By applying $F_*$ to the previous 
            equation we get $[u_M,v_M]=-[u,v]_M$.
            \item[3.] We have $g_*(u_M)=\Ad_g(u)_M$. This is true in $\overline{M}=G\times M$. In fact, we have 
            \begin{equation}
                \begin{split}
                    \pi_G*(a,x)g_{*,(a,x)}u_{\overline{M}}(a,x)
                    &=(\pi_G\circ\psi_g)_{*,(a,x)}u_{\overline{M}}(a,x)\\
                    &=\frac{d}{dt}\bigg|_{t=0}\pi_G\circ\psi_g\circ(r_{\exp(-tu)}(a),x)\\
                    &=\frac{d}{dt}\bigg|_{t=0}\pi_G(r_{g^{-1}}r_{\exp(-tu)}(a),x)\\
                    &=\frac{d}{dt}\bigg|_{t=0}a\exp(-tu)g^{-1}\\
                    &=\frac{d}{dt}\bigg|_{t=0}ag^{-1}c_g(\exp(-tu))\\
                    &=\frac{d}{dt}\bigg|_{t=0}ag^{-1}\exp(-\Ad_g(u))\\
                    &=\frac{d}{dt}\bigg|_{t=0}r_{\exp(-\Ad_g(u))}(ag^{-1})\\
                    &=-(\Ad_g(u))^L(ag^{-1}).
                \end{split}
            \end{equation}
            Finally 
            \begin{equation}
                \begin{split}
                    g_*u_{\overline{M}}\big|_{g\cdot(a,x)}&=-(\Ad_g(u)^L(ag^{-1}),0)\\
                    &=-(\Ad_g(u)^L,0)\big|_{g\cdot(a,x)}.
                \end{split}
            \end{equation}
        \end{itemize}
        Now we have 
        \begin{equation}
            g_*u_M=g_*F_*u_{\overline{M}}=F_*g_*u_{\overline{M}}=F_*(\Ad_g(u))_{\overline{M}}=(\Ad_g(u))_M.
        \end{equation}
        \item Let $(M,G,\psi)$ be a $G$-manifold. For $x\in M$, we define the isotropy subgroup 
        $G_x=\{g\in G:gx=x\}$. It's clear that $G_x\subset G$ is closed, so it's a Lie subgroup of $G$. Define 
        $\fr{g}_x=\{u\in\fr{g}:u_M(x)=0\}$, we claim that  
        \begin{equation}
            \text{Lie}(G_x)=\fr{g}_x.
        \end{equation}
        Let $\exp(tu)\in G_x,\;\forall t\in\bbR$, that is, assume that $u\in\text{Lie}(G_x)$. Then, by definition, we have 
        $\exp(tu)x=x$. So, by differentiating this equation we get $u_M(x)=0$ and $u\in\fr{g}_x$. Conversely, assume that 
        $u\in\fr{g}_x$. Let $g=\exp(ts),\;s\in\bbR$, then 
        \begin{equation}
            \frac{d}{dt}\bigg|_{t=s}\exp(tu)x= g_{*,x}(\frac{d}{dt}\bigg|_{t=0}\exp(tu)x)=0
        \end{equation}
        and $\exp(tu)\cdot x = x,\forall t\in\bbR$. Therefore $u\in\text{Lie}(G_x)$.
        \item Now, if $\mathcal{O}_x=Gx$, it's a fact that if we identify $G/G_x\cong \mathcal{O}_x$ by the 
        orbit stabilizer theoerem, then $\mathcal{O}_x$ inherits a differentiable structure such that the inclusion 
        $i:\mathcal{O}_x\hookrightarrow M$ is an imersion. In this context we claim that 
        \begin{equation}
            i_{*,x}T_x\mathcal{O}_x=\{u_M(x): u\in\fr{g}\}.
        \end{equation}
        Lets prove it. We know that $i:G/G_x\hookrightarrow M:[g]\mapsto gx$ is an imersion. First observe that 
        \begin{equation}
            \begin{split}
                \frac{d}{dt}\bigg|_{t=0}i\circ\pi\circ r_{\exp(tu)}(g)&=\frac{d}{dt}\bigg|_{t=0}g\exp(tu)\cdot x\\
                &=g_{*,x}u_M(x)\\
                &=(\Ad_g(u))_M(gx).
            \end{split}
        \end{equation}
        Since $\pi:G\rightarrow G/G_x$ is a submersion, we conclude that 
        $i_{*,[g]}T_{[g]}(G/G_x)\subset \{u_M(g\cdot x):u\in\fr{g}\}$. Now define the map 
        \begin{equation}
            \fr{g}\to \{u_M(g\cdot x):u\in\fr{g}\}
        \end{equation}
        which has kernel equal to $\text{Lie}(G_{g\cdot x})$. Therefore we have an isomorphism
        \begin{equation}
            \fr{g}/\text{Lie}(G_{g\cdot x})\overset{\cong}{\to}\{u_M(g\cdot x):u\in\fr{g}\}.
        \end{equation}
        Therefore the dimension of the last subspace is equal to 
        \begin{equation}
            \begin{split}
                \dim(G)-\dim(G_{g\cdot x})&=\dim(G)-\dim(gG_xg^{-1})\\
                &=\dim(G)-\dim(G_x)\\
                &=\dim(G/G_x).
            \end{split}
        \end{equation}
        Therefore, by dimensional reasons we have $i_{*,[g]}T_{[g]}(G/G_x)=\{u_M(g\cdot x):u\in\fr{g}\}$ and since 
        $i_*$ is injective we can identify $T_x{\mathcal{O}_x}=\{u_M(x):u\in\fr{g}\}$.
\end{itemize}

\subsection{Co-adjoint orbits and KKS form}

    Let $G$ be a Lie group. We will show that the orbits on $\fr{g}^*$ under the co-adjoint action have a natural
    structure of a symplectic manifold.\par
    \textbf{KKS}: Let $\xi\in\mathcal{O}$, we define $\omega\in\bigwedge^2T_\xi^*\mathcal{O}$ by the formula: 
    \begin{equation}
        \omega_\xi(u_{\fr{g}^*}(\xi),v_{\fr{g}^*}(\xi))=\langle\xi,[u,v] \rangle.
    \end{equation}
    We need to verify a few things:
\begin{itemize}
    \item $\omega_\xi$ is well defined. Assume that $u_{\fr{g}^*}(\xi)=u_{\fr{g}^*}^\prime(\xi)$ and 
    $v_{\fr{g}^*}(\xi)=v_{\fr{g}^*}^\prime(\xi)$, then 
    \begin{equation}
        \begin{split}
            \omega_\xi(u_{\fr{g}^*}^\prime(\xi),v_{\fr{g}^*}^\prime(\xi))&=
            -u_{\fr{g}^*}^\prime(\xi)(v^\prime)\\
            &=-u_{\fr{g}^*}(\xi)(\xi)(v^\prime)\\
            &=\langle \xi,[u,v^\prime] \rangle\\
            &=v_{\fr{g}^*}^\prime(\xi)(u)\\
            &=v_{\fr{g}^*}(\xi)(u)\\
            &=\langle\xi,[u,v] \rangle.
        \end{split}
    \end{equation}
    \item It's skew symmetric and nondegenerate (clear): $\langle\xi,[u,v] \rangle=0\;\forall v\iff u_{\fr{g}^*}(\xi)=0$.
    \item It's $G-$invariant, so it's smooth. This can be verified directly by using the relation $g_*(u_{\fr{g}^*})=(\Ad_g(u))_{\fr{g}^*}$.
    \item By the previous item we see that $\psi_{e^{tu}}^*\omega=\omega$. Therefore 
    \begin{equation}
        \mathcal{L}_{\fr{g}^*}\omega=0\iff (d\iota_{u_{\fr{g}^*}}+\iota_{u_{\fr{g}^*}}d)\omega=0.
    \end{equation}
    \textbf{Claim}: We have $\iota_{u_{\fr{g}^*}}\omega=du$ where we identify $u\in\fr{g}\cong C^\infty_{\text{Lin}}(\fr{g}^*)$.\par 
    This can be seen directly by the equation:
    \begin{equation}
        \begin{split}
            \iota_{u_{\fr{g}^*}}\omega\big|_\xi(v_{\fr{g}^*}(\xi))&=\langle\xi,[u,v] \rangle\\
            &=\langle\xi,[u,v] \rangle\\
            &=v_{\fr{g}^*}(\xi)(u)\\
            &=du\big|_\xi v_{\fr{g}^*}(\xi)\;(\text{by the identification }\fr{g}^{**}=\fr{g})
        \end{split}
    \end{equation}
    \item For $G=SO(3)$ the co-adjoint and adjoint orbits are the same. We have an isomorphism of Lie algebras
    $\bbR^3\cong\fr{so}(3)$ which maps $e_i\in\bbR^3$ to the associated matrix representation of the map 
    $(e_i)_\times:v\mapsto e_i\times v$. Moreover, we have the identity 
    \begin{equation}
        A(u\times v)= Au\times Av,\; \forall A\in SO(3).
    \end{equation}
    With this equation we get $(Au)_\times v = Au\times v = Au_\times A^{-1}(v)=\Ad_A(u_\times)v$. So this isomorphism 
    is $SO(3)-$equivariant, with the adjoint orbits being spheres in $\bbR^3$.
    \begin{dang}
     \begin{remark}
        If $\nu: \fr{g}\to \fr{g}^*$ is an isomorphism that intetwine the adjoint and co-adjoint actions of $G$. Then we 
        can induce a symplectic form on adjoint orbits using KKS. The formula is the one 
        \begin{equation}
            \omega_\xi(u_{\fr{g}}(\xi),v_{\fr{g}}(\xi))=\langle \nu(\xi), [u,v]\rangle.
        \end{equation}
    \end{remark}    
    \end{dang}

    Going back to definitions: If $\xi\in(\bbR^3)^*$, we have 
    \begin{equation}
        \begin{split}
            \omega_\xi(u_{\fr{g}}(\xi),v_{\fr{g}}(\xi))&=\langle\xi,[u,v]\rangle\\
        &=\langle \xi, u\times v\rangle \;(\bbR^3\text{ is a Lie algebra with the croos product})\\
        &=\xi\cdot (u\times v)\;(\text{By the isomorphism }\bbR^3\cong(\bbR^3)^*)
        \end{split}
    \end{equation}
    \begin{dang}
    \begin{remark}
        Let $S\subset M$ a hypersurface, $M$ oriented with orientation form $\omega$. Let 
        $N\in\fr{X}(M)$ be a vector field nowhere tangent to $S$, then $\iota_N\omega|_S$ is an orientation 
        form for $S$.
    \end{remark}        
    \end{dang}

    By the previous remark we can use the vector field 
    $N= x\tfrac{\partial}{\partial x}+y\tfrac{\partial}{\partial y}+z\tfrac{\partial}{\partial z}$ to induce 
    an orientation form 
    \begin{equation}
        \sigma = xdy\wedge dz - ydx\wedge dz+zdx\wedge dy \big|_{\mathbb{S}^2}
    \end{equation}
    on the spheres. In fact, if $f:(x,y,z)\mapsto x^2+y^2+z^2-1$, then $df(N)=2(x^2+y^2+z^2)$ is never zero on $\mathbb{S}^2$, so 
    $N$ is nowhere tangent to $\mathbb{S}^2$.\par 

    \begin{dang}
        \begin{remark}
            If $f:\bbR^n\to\bbR$ is smooth with regular value $c$, then $S=f^{-1}(0)$ is a regular submanifold 
            with tangent spaces 
            \begin{equation}
                i_{*p}T_pS \cong \nabla f(p)^\perp
            \end{equation}
        \end{remark}
    \end{dang}

    In our case we have $\nabla f(\xi) = \xi$, so $T_\xi\mathbb{S}^2=\xi^\perp=\{\xi\times u:u\in\bbR^3\}$. If we 
    use the area form 
    \begin{equation}
        \begin{split}
            \sigma_\xi(\xi\times u,\xi\times v) &=\xi\cdot ((\xi\times u)\times (\xi\times v))\;(\text{use the area form})\\
            &=-(\xi\times u)\cdot (\xi\times(\xi\times v))\\
            &=-(\xi\times u)\cdot (\xi(\xi\cdot v)-v(\xi\cdot\xi))\\
            &=\|\xi\| (\xi\times u)\cdot v \\
            &=\xi\cdot (u\times v)\|\xi\|.
        \end{split}
    \end{equation}

    In conclusion, we see the relation between the KKS and the area form:
    \begin{equation}
        \omega_{\text{kks}}=\frac{1}{r}\sigma
    \end{equation}
    where $r=\|\xi\|$ is the radius of the sphere.
\end{itemize}

\subsection{Hamiltonian actions}
Let $X\in T_0\bbR^n\cong \bbR^n$, then $\widetilde{X}(g)=(g,X)$. In fact
\begin{equation}
    \begin{split}
         \ell_{g*,0}(\frac{\partial}{\partial x_i}\bigg|_g)x_j &= \frac{\partial}{\partial x_i}\bigg|_g x_j\circ \ell_g\\
         &=\frac{\partial}{\partial x_i}\bigg|_g(x_j+x_j(g))=\delta_{ij}.
    \end{split}
\end{equation}

In general if $\widetilde{X}\in\fr{X}_{\text{inv}}(G)$, then we set $\exp(X)=\gamma_1^{\widetilde{X}}(e)$. If $G=\bbR^n$, then 
$\gamma_t^{\widetilde{X}}(0)=tX$ and $e^X=X$.


    Let $(M,\omega)$ be a symplectic manifold. Then

    \begin{itemize}
        \item $G-$action is symplectic if $\mathcal{L}_X\omega = 0$. This condition is equivalent to 
        $\psi_g^*\omega=\omega\;\forall g\in G$ if $G$ is connected.
        \item $G-$action is Hamiltonian if there exist a linear map $\hat{\mu}:\fr{g}\to C^\infty(M)$
        such that the following diagram commutes: 
        \begin{center}
            \begin{tikzcd}
                                             &   & C^\infty(M)\arrow[dd, dashed]& \hat{\mu}(u)\arrow[dd, mapsto] \\
    \fr{g}\arrow[rru, "\hat{\mu}"]\arrow[rrd]&u\arrow[rru, mapsto]\arrow[rrd, mapsto]  &    &\\
                                             &   & \fr{X}(M)   &u_M = X_{\hat{\mu(u)}}
            \end{tikzcd}
        \end{center}
    \end{itemize}

    \begin{remark}
        We have a bijection 
        \begin{equation}
            \{\hat{\mu}: \fr{g}\to C^\infty(M)\} \leftrightarrow \{\mu:M\to\fr{g}^*\}
        \end{equation}
    \end{remark}

\begin{itemize}
    \item The action is weakly Hamiltonian if
    \begin{equation}
        u_M = X_{\hat{\mu}(u)} \iff \iota_{u_M}\omega = d\hat{\mu} = d\langle\mu,u\rangle.
    \end{equation}
    \item The action is Hamiltonian if, in addition, we have 
    \begin{equation}
        \mu \circ \psi_g = \Ad_g^*\circ\mu \overset{G\text{ connected}}{\iff} \hat{\mu}([u,v])=\{\hat{\mu}(u),\hat{\mu}(v)\}.
    \end{equation}
    The relation between the two definitions is $\hat{\mu}(u)\big|_x = \langle\mu(x),u\rangle$.
\end{itemize}

Now, look at the following equation for $x\in M, v\in \fr{g}$:
\begin{equation}
    \langle \mu\circ\psi_g(x),v\rangle = \langle\Ad_g^*\circ\mu(x),v\rangle = \langle \mu(x), \Ad_{g^{-1}}(v)\rangle.
\end{equation}
This equation means that $\hat{\mu}(v)\circ\psi_g(x)=\hat{\mu}(\Ad_{g^{-1}}(v))\big|_x$. If we put $g=e^{tu},u\in\fr{g}$ and 
differentiate this equation, we get $\hat{\mu}([u,v])=\{\hat{\mu}(u),\hat{\mu}(v)\}$. For the converse we need the hypothesis 
of $G$ being connected.\par 

\begin{remark}
    Remember that $\fr{X}_{\text{Ham}}(M)\subset\fr{X}_{\text{Sym}}(M)$. 
\end{remark}

In our case we have $\mathcal{L}_{u_M}\omega = 0$. Since $u_M$ has flux given by $\psi_{e^{tu}}$, we get 
$\psi_{e^{tu}}^*\omega=\omega$. Therefore $\psi_g^*\omega = \omega,\;\forall g\in G^\circ$.
    
\begin{dang}
    \begin{theorem}
        (Noether principle). Let $(M,\omega,\mu)$ be a Hamiltonian $G$-space and $H\in C^\infty(M)$. Then 
        $H$ is $G$-invariant if and only if $\mu$ is preserved by the Hamiltonian flux of $H$.
    \end{theorem}
\end{dang}
Proof: Remember that $H$ is $G$ invariant if $\psi_g^*H=H\iff\mathcal{L}_{u_M}H=0$ if $G$ is connected. Also, 
$\mu$ is preserved by the Hamiltonian flux of $H$ if $\varphi_t^*\mu=\mu$ where $\varphi_t=\varphi_t^{X_H}$ is the 
flux of $X_H$. This equation is true if 
\begin{equation}
    \begin{split}
        \mathcal{L}_{X_H}\langle\mu,u\rangle &= \frac{d}{dt}\bigg|_{t=0}\varphi_t^*\langle\mu,u\rangle\\
           &=\frac{d}{dt}\bigg|_{t=0}\langle\varphi_t^*\mu,u\rangle=0\;\forall u\in\fr{g}\\
           &\iff \varphi_t^*\mu =\mu,\;\forall t\in\bbR.
    \end{split}
\end{equation}
But we see that 
\begin{equation}
    \begin{split}
        \mathcal{L}_{X_H}\langle\mu,u\rangle &= (d\iota_{X_H}+\iota_{X_H}d)\langle\mu,u\rangle\\
        &=\iota_{X_H}d\langle\mu,u\rangle\\
        &=\iota_{X_H}\iota_{u_M}\omega\\
        &=-dH(u_M)\\
        &=-\mathcal{L}_{u_M}H.
    \end{split}
\end{equation}
\qed

\begin{example}
    A general situation: $(M,\omega)$ is symplectic with $\omega = -d\theta$ exact. Suppose that we have 
    a $G$ on $M$ such that $\psi_g^*\omega = \omega$. In this case we define 
    \begin{equation}
        \iota_{u_M}\theta =  \langle \mu,u\rangle.
    \end{equation}
    It's easy to see that this defines a Hamiltonian action in $M$.
\end{example}

\begin{example}
    (Cotangent lifting). Take a $G$-action $\psi_g$ on $M$, then we can lift 
    it to a $G$-action $\widehat{\psi}_g$ on $T^*M$. So we have $\widehat{\psi}_g^*\alpha=\alpha$ and thus we get 
    a Hamiltonian map given by $\langle\mu,u\rangle = \iota_{u_M}\alpha$. 
\end{example}

\begin{example}
    Let $G=\bbR^n$ acting on $M=\bbR^n$ by translation $g\cdot q = q+g$. With the identification 
    $T^*M\cong\bbR^n\times\bbR^n$ we have $g\cdot(q,p)=(g+q,p)$. We can explictly compute the moment map:
    \begin{equation}
        \begin{aligned}
            &u = (a,b) \mapsto u_M = \sum_ia_i\frac{\partial}{\partial q_i}+\sum_ib_i\frac{\partial}{\partial p_i}\\
            &\iota_{u_M}\sum_ip_idq_i=\sum_ip_idq_i(u_M)=\sum_ip_ia_i\\
            &\mu(q,p):\bbR^n\to\bbR^n:e_i\mapsto p_i 
        \end{aligned}
    \end{equation}
    So $\mu(q,p)=p$ in the choosen basis.
\end{example}

\begin{example}
    Let $G=SO(3)$ acting in $\bbR^3$ in the usual way a lift it to an action on $T^*\bbR^3=\bbR^3\times\bbR^3$. 
    Now observe that 
    \begin{equation}
    \begin{split}
    \psi_{g*,v}(\frac{\partial}{\partial x_i}\bigg|_v)x_j&=\frac{\partial}{\partial x_i}\bigg|_vx_j\circ \psi_g \\
    &=d(x_j\circ \psi_g)_v(e_i)\\
    &=x_j \circ\psi_g(e_i)\\
    &=g_{ji}
    \end{split}
    \end{equation}
    This proves that $\psi_{g*,v}(\sum_ia_i\frac{\partial}{\partial x_i}\bigg|_v)
    =\sum_i(ga)_i\frac{\partial}{\partial x_i}\bigg|_{gv}$. Therefore, the corresponding action on $T\bbR^3$ is 
    \begin{equation}
        g\cdot(q,p)=(gq,gp).
    \end{equation}
    Remember that the inverse of the dualization preserve base points and inverse the dual map, so the action on 
    $T^*\bbR^3$ is 
    \begin{equation}
        g\cdot (q,p)=(gq,(g^{-1})^Tp)=(gq,gp),\;g\in SO(3).
    \end{equation}
    Finnaly we see that 
    \begin{equation}
        \begin{aligned}
            &u_M(q,p)=\frac{d}{dt}\bigg|_{t=0}(e^{tu}q,e^{tu}p)=(uq,up)\\
            &u_M = \sum_i (uq)_i\frac{\partial}{\partial q_i}+\sum_i (up)_i\frac{\partial}{\partial p_i}\\
            &\iota_{u_M}\sum_ip_idq_i = \sum_i p_iu_{ij}q_j\\
            &\mu(q,p):\fr{so}(3)\to \bbR: \ell_1\mapsto p_3q_2-p_2q_3,\dots\Rightarrow \mu(q,p)=q\times p
        \end{aligned}
    \end{equation}
\end{example}
\subsection{Marsden, Weinstein, Meyer theorem}